\section{What Preservation Can Certify}
\label{sec:ranktheory}
Distillation controls a probability distribution, whereas retrieval depends on the highest-ranked relevant candidate.
The following result gives the exact connection on a fixed gallery.
Let $q$ and $p$ denote teacher and student probabilities on that gallery.
Let $R$ contain its relevant candidates.
Define $C_K(q,R)$ as the smallest forward KL divergence needed to place $K$ irrelevant candidates above every relevant candidate.
Ties count as failures.

\begin{proposition}
Assume the teacher is strictly correct and at least $K$ irrelevant candidates exist.
Select the $K$ largest irrelevant teacher probabilities.
Cap relevant probabilities above a level $t$ and raise selected irrelevant probabilities below $t$.
Choose $t$ so the removed and added mass agree.
Call the resulting distribution $p^*$.
Then $C_K(q,R)=\KL(q\|p^*)$.
Consequently, $\KL(q\|p)<C_K(q,R)$ guarantees a relevant candidate within the first $K$ positions.
\end{proposition}

Appendix~\ref{app:ranktheory} gives the formula, attaining projection, proof, and numerical boundary cases.
The radius treats all valid candidates jointly.
It can exceed a bound that protects only the teacher's best relevant candidate.
The same-gallery requirement matters: a minibatch penalty cannot certify unseen distractors in a larger gallery.
The empirical practical criterion therefore uses measured full-pool retrieval.
